\chapter{Introduction}
\label{cha:introduction}

\section{Purpose of this document}
\label{sec:purpose}

This manual documents \texttt{DEMcookies}, a compact two- and three-dimensional Discrete
Element Model (DEM) written in Python.  The code simulates an assembly of
circular, polydisperse particles that are first deposited into a rectangular
container under gravity (\emph{raining} deposition) and subsequently compressed
by a rigid loading platen that transmits a prescribed, time-dependent force.
Most of this manual describes the planar model, which Sections~\ref{sec:3d} and~\ref{sec:3d-numerics} extend to three dimensions.

The document is intended to be self-contained: every equation that the code
evaluates is stated explicitly, every numerical parameter is defined, and every
algorithmic choice is justified.  The reader should be able to (i) reproduce the
implementation from the equations alone, (ii) understand the meaning and the
admissible range of each input parameter, and (iii) judge the accuracy and the
limitations of the results.

The manual is organised as follows.
\begin{itemize}[leftmargin=1.6em]
  \item Chapter~\ref{cha:introduction} introduces the discrete element method,
        the physical problem, the notation and the structure of the package.
  \item Chapter~\ref{cha:formulation} develops the complete mathematical
        formulation: kinematics, the Newton--Euler equations of motion, the
        Hertz--Mindlin contact law, boundary conditions and the platen
        sub-system.
  \item Chapter~\ref{cha:numerical} describes the numerical solution: time
        integration, stability, contact detection, history bookkeeping,
        deposition algorithm and computational complexity.
  \item Chapter~\ref{cha:examples} presents verification tests with analytical
        reference solutions and a full application example.
\end{itemize}

\section{The discrete element method}
\label{sec:dem}

The discrete element method, introduced by Cundall and Strack for granular
assemblies~\cite{cundall1979}, treats a granular material as a large collection
of individual rigid bodies whose motion is governed by Newton's second law.  The
particles interact only through contacts.  Two families of formulations exist:

\begin{description}[leftmargin=1.6em,style=nextline]
  \item[Hard-sphere (event-driven) methods] Particles are perfectly rigid,
        collisions are instantaneous, and the simulation advances from collision
        event to collision event.  These methods are efficient for dilute,
        rapidly flowing systems but cannot represent enduring, multi-body
        contacts.
  \item[Soft-sphere (time-driven) methods] Particles are allowed to overlap
        slightly; the overlap is interpreted as the local elastic deformation of
        the contact region, and a constitutive law converts it into a contact
        force.  The equations of motion are integrated with a small, fixed time
        step.
\end{description}

\texttt{DEMcookies} is a soft-sphere code.  The soft-sphere paradigm is the
standard choice for dense, quasi-static problems such as the confined
compression considered here, because it naturally handles persistent contacts,
force chains and static equilibrium~\cite{osullivan2011,zhu2007}.

The soft-sphere DEM cycle repeated at every time step is:
\begin{enumerate}[leftmargin=2.2em]
  \item \textbf{Contact detection.}  Find all pairs of bodies whose surfaces
        overlap (broad phase followed by a narrow, exact geometric test).
  \item \textbf{Force evaluation.}  For each contact, compute the normal and
        tangential forces from the overlap, the relative surface velocity and
        the accumulated tangential displacement; add gravity.
  \item \textbf{Integration.}  Advance translational and rotational degrees of
        freedom with an explicit integrator.
\end{enumerate}

\section{The physical problem}
\label{sec:problem}

\begin{figure}[htbp]
  \centering
  \setlength{\unitlength}{1mm}
  \begin{picture}(90,62)
    % container
    \put(25,5){\line(1,0){40}}          % bottom wall
    \put(25,5){\line(0,1){50}}          % left wall
    \put(65,5){\line(0,1){50}}          % right wall
    % platen
    \linethickness{1.2pt}
    \put(25,40){\line(1,0){40}}
    \linethickness{0.4pt}
    % load arrow
    \put(45,52){\vector(0,-1){11}}
    \put(47,47){$F(t)$}
    % particles (schematic)
    \put(30,9){\circle{7}}   \put(38,8){\circle{5}}   \put(45,10){\circle{8}}
    \put(54,8){\circle{5}}   \put(61,9){\circle{6}}
    \put(31,16){\circle{6}}  \put(39,15){\circle{7}}  \put(48,17){\circle{5}}
    \put(56,16){\circle{7}}  \put(62,15){\circle{4}}
    \put(29,23){\circle{5}}  \put(36,23){\circle{6}}  \put(44,24){\circle{7}}
    \put(52,23){\circle{5}}  \put(60,23){\circle{7}}
    \put(31,30){\circle{7}}  \put(40,31){\circle{6}}  \put(48,30){\circle{7}}
    \put(57,31){\circle{6}}  \put(63,30){\circle{4}}
    \put(30,36){\circle{5}}  \put(38,37){\circle{6}}  \put(46,36){\circle{5}}
    \put(54,37){\circle{6}}  \put(61,36){\circle{5}}
    % dimensions
    \put(25,60){\vector(1,0){40}}
    \put(65,60){\vector(-1,0){40}}
    \put(43,62){$B$}
    \put(20,5){\vector(0,1){50}}
    \put(20,55){\vector(0,-1){50}}
    \put(14,29){$L$}
    % axes
    \put(8,5){\vector(1,0){9}}  \put(18,4){$x$}
    \put(8,5){\vector(0,1){9}}  \put(6,15){$y$}
    % gravity
    \put(76,45){\vector(0,-1){12}} \put(78,38){$g$}
  \end{picture}
  \caption{Model problem: a rectangular container of width $B$ and height $L$
  with fixed left, right and bottom walls.  Polydisperse circular particles are
  rained in under gravity; a rigid platen then closes the top and transmits the
  prescribed load $F(t)$.}
  \label{fig:problem}
\end{figure}

The reference problem solved by the package (Figure~\ref{fig:problem}) consists
of two consecutive stages.

\paragraph{Stage 1 -- gravitational deposition.}
An initially empty box of width $B$ and height $L$, with fixed walls at $x=0$,
$x=B$ and $y=0$, is filled by repeatedly inserting small batches of
non-overlapping particles in a thin band just below the top of the domain.  The
particles fall under gravity, collide, dissipate energy and build a dense random
packing.  Insertion continues until the growing bed reaches the insertion band,
after which the assembly is allowed to relax until its mean kinetic energy falls
below a quiescence tolerance.  This procedure is the numerical analogue of
pluviation (sand raining) used in laboratory sample preparation.

\paragraph{Stage 2 -- confined compression.}
A rigid, massive platen spanning the full width $B$ is placed on top of the
settled bed.  The platen is constrained to move vertically only and obeys its
own equation of motion, driven by the reaction of the particles below and by an
externally prescribed load $F(t)$.  Because the lateral walls are rigid and
fixed, the test is an \emph{oedometric} (one-dimensional, laterally confined)
compression.  The recorded platen displacement and reaction force yield the
macroscopic load--displacement response of the packing.

\paragraph{Other problems.}
The rectangular box is one instance of a more general set-up.  The domain may be
any simple polygon, each of its edges is a fixed wall (rough or smooth), an open
boundary through which particles can leave, or a loaded platen, and every
platen follows its own load history: a pulse, a Heaviside step (optionally
ramped), a sinusoid or a user-defined time history.  The location at which a
load is applied is set by the geometry, so a strip footing is modelled by
splitting the top side into collinear edges and loading only the middle one
(Section~\ref{sec:walls}).  A platen may also span several edges with different
orientations, which then move as one rigid body, so a load can be applied
through a wedge or a curved indenter rather than along a straight line
(Section~\ref{sec:platen}).  Finally, the particles can be immersed in water
below a water table, which may rise or fall with time, and an upward or
downward seepage flow can be imposed through them.  The water acts on the
particles through buoyancy and porosity-dependent drag, and, if requested, an
excess pore pressure that builds up under fast loading and dissipates by flow,
so that undrained loading and consolidation can be studied
(Section~\ref{sec:water}).

\section{Structure of the package}
\label{sec:structure}

The package is deliberately small; each module maps onto one section of this
manual (Table~\ref{tab:modules}).

\begin{table}[htbp]
  \centering
  \small
  \caption{Modules of \texttt{DEMcookies} and the corresponding sections.}
  \label{tab:modules}
  \begin{tabular}{l>{\raggedright\arraybackslash}p{4.9cm}>{\raggedright\arraybackslash}p{4.3cm}}
    \toprule
    Module & Main contents & Documented in \\
    \midrule
    \texttt{contact.py} &
      \texttt{MaterialProperties}, \texttt{HertzMindlinModel},
      \texttt{rayleigh\_time\_step} &
      Sections~\ref{sec:normal}--\ref{sec:tangential},
      Section~\ref{sec:timestep} \\
    \texttt{geometry.py} &
      \texttt{Geometry} (polygonal domain) &
      Section~\ref{sec:domain} \\
    \texttt{geometry3d.py} &
      \texttt{Geometry3D} (polyhedral domain) &
      Section~\ref{sec:3d} \\
    \texttt{facets.py} &
      closest points, contacts and seating of edges and faces &
      Sections~\ref{sec:coupling} and~\ref{sec:3d-numerics} \\
    \texttt{vectors.py} &
      vector operations shared by 2D and 3D &
      Section~\ref{sec:3d-numerics} \\
    \texttt{boundary.py}, \texttt{segments.py} &
      \texttt{Wall}, \texttt{Open}, \texttt{LoadedPlaten}, \texttt{Platen} &
      Sections~\ref{sec:walls}--\ref{sec:platen} \\
    \texttt{loads.py} &
      \texttt{Pulse}, \texttt{Heaviside}, \texttt{Sinusoidal},
      \texttt{TimeHistory} &
      Section~\ref{sec:platen} \\
    \texttt{neighbours.py} &
      cell grid, \texttt{NeighbourList} &
      Section~\ref{sec:broadphase} \\
    \texttt{water.py} &
      \texttt{Water}, buoyancy, drag, porosity grid, pore pressure &
      Sections~\ref{sec:water} and~\ref{sec:water-numerics} \\
    \texttt{particle\_generation.py} &
      \texttt{Particle}, \texttt{Particle3D}, \texttt{generate\_lattice\_particles} &
      Section~\ref{sec:deposition} \\
    \texttt{solver.py} &
      \texttt{VelocityVerlet}, \texttt{DEMSimulation}, \texttt{RainSettings} &
      Chapter~\ref{cha:numerical} \\
    \texttt{globals.py} &
      constants, \texttt{GlobalSettings} &
      Chapter~\ref{cha:numerical} \\
    \texttt{plotting.py} &
      animation, load curves &
      Section~\ref{sec:recording} \\
    \texttt{vtk\_export.py} &
      \texttt{VTKWriter}, export to ParaView &
      Section~\ref{sec:recording} \\
    \texttt{main.py}, \texttt{examples/} &
      driver scripts &
      Chapter~\ref{cha:examples} \\
    \texttt{tests/} &
      unit and verification tests &
      Section~\ref{sec:unit-tests} \\
    \texttt{manual/scripts/} &
      scripts that produce the results of this manual &
      Chapter~\ref{cha:examples} \\
    \bottomrule
  \end{tabular}
\end{table}

\section{Modelling assumptions}
\label{sec:assumptions}

The following assumptions are built into the formulation and should be kept in
mind when interpreting results.

\begin{enumerate}[leftmargin=2.2em,label=A\arabic*.]
  \item \textbf{Two-dimensional kinematics.}  Motion is restricted to the
        $x$--$y$ plane.  Each particle has three degrees of freedom: two
        translations and one rotation about the out-of-plane axis.
        In a polyhedral domain the particles move in space instead, with six degrees of freedom (Section~\ref{sec:3d}).
  \item \textbf{Spherical mass properties.}  Although the geometry is a disc,
        the mass and the moment of inertia are those of a \emph{sphere},
        $m=\tfrac{4}{3}\pi\rho R^{3}$ and $I=\tfrac{2}{5}mR^{2}$.  The model is
        therefore best interpreted as a planar arrangement of spheres rather
        than as a plane-strain slice of a cylinder
        (see Section~\ref{sec:massprops}).
  \item \textbf{Single material.}  All particles and all boundaries share the
        same density, Young's modulus, Poisson ratio and restitution
        coefficient, so the effective contact constants $E^{*}$, $G^{*}$ and
        $\beta$ are evaluated once at construction time.  The friction
        coefficient is that of the material, but may be overridden for each
        wall or platen.
  \item \textbf{Small overlaps.}  The Hertz--Mindlin law is a small-strain
        elastic solution.  It is quantitatively valid only while the overlap
        remains a small fraction of the particle radius
        (Section~\ref{sec:limitations}).
  \item \textbf{No cohesion.}  Contacts transmit only compression, viscous
        damping, Coulomb friction and, optionally, a rolling resistance moment
        that stands in for the angular shape of real
        grains~(Section~\ref{sec:rolling}).  Adhesive forces, liquid bridges and
        particle breakage are not modelled.
  \item \textbf{Dry or saturated system.}  By default there is no interstitial
        fluid.  Optionally the particles are immersed in water below a water
        table.  The water is either drained -- its pressure is hydrostatic and
        any flow is prescribed -- or coupled, with an excess pore pressure that
        follows a Biot storage equation with Darcy flow on a coarse grid.  The
        fluid itself is not resolved: it has no inertia and no viscous stresses
        of its own (Section~\ref{sec:water}).
  \item \textbf{Rigid boundaries.}  Walls are perfectly rigid, straight
        segments (planar faces in 3D) of infinite mass and infinite radius; a platen is rigid but has
        a finite mass and one translational degree of freedom, along its
        direction of motion.
\end{enumerate}

\section{Notation}
\label{sec:notation}

Vectors are set in bold, $\mathbf{a}=(a_{x},a_{y})$.  Because the model is
planar, angular quantities are scalars measuring the out-of-plane
($z$) component, and the cross product of two in-plane vectors is the scalar
$(\mathbf{a}\times\mathbf{b})_{z}=a_{x}b_{y}-a_{y}b_{x}$.
In three dimensions (Section~\ref{sec:3d}) vectors have a third component, $a_{z}$, and angular quantities and cross products are vectors.
Table~\ref{tab:notation} collects the symbols used throughout the manual.

\begin{table}[htbp]
  \centering
  \small
  \caption{Notation.}
  \label{tab:notation}
  \begin{tabular}{llll}
    \toprule
    Symbol & Meaning & Unit & Code name \\
    \midrule
    $N$              & number of particles              & --                   & \texttt{sim.n} \\
    $B,\ L$          & box width and height             & \unit{\metre}        & \texttt{Geometry.rectangle} \\
    $\mathbf{p}_{k}$, $\mathbf{n}_{k}$ & vertex, inward edge normal & \unit{\metre}, -- & \texttt{vertices}, \texttt{inward\_normals} \\
    $\mathbf{x}_{i}$ & particle centre                  & \unit{\metre}        & \texttt{pos[i]} \\
    $\mathbf{v}_{i}$ & particle velocity                & \unit{\metre\per\second} & \texttt{vel[i]} \\
    $\omega_{i}$     & angular velocity ($z$)           & \unit{\radian\per\second} & \texttt{omega[i]} \\
    $R_{i}$          & particle radius                  & \unit{\metre}        & \texttt{radius[i]} \\
    $m_{i}$          & particle mass                    & \unit{\kilogram}     & \texttt{mass[i]} \\
    $I_{i}$          & moment of inertia                & \unit{\kilogram\metre\squared} & \texttt{inertia[i]} \\
    $\rho$           & material density                 & \unit{\kilogram\per\cubic\metre} & \texttt{density} \\
    $E$              & Young's modulus                  & \unit{\pascal}       & \texttt{youngs\_modulus} \\
    $\nu$            & Poisson ratio                    & --                   & \texttt{poisson\_ratio} \\
    $G$              & shear modulus                    & \unit{\pascal}       & \texttt{shear\_modulus} \\
    $E^{*},\ G^{*}$  & effective contact moduli         & \unit{\pascal}       & \texttt{e\_star}, \texttt{g\_star} \\
    $e$              & coefficient of restitution       & --                   & \texttt{restitution} \\
    $\beta$          & damping ratio derived from $e$   & --                   & \texttt{beta} \\
    $\mu$            & Coulomb friction coefficient     & --                   & \texttt{friction} \\
    $g$              & gravitational acceleration       & \unit{\metre\per\second\squared} & \texttt{gravity} \\
    $\delta_{n}$     & normal overlap                   & \unit{\metre}        & \texttt{overlap} \\
    $\mathbf{n},\ \mathbf{t}$ & contact normal and tangent & --                 & \texttt{normal}, \texttt{tangent} \\
    $R^{*},\ m^{*}$  & effective radius and mass        & \unit{\metre}, \unit{\kilogram} & \texttt{effective\_radius}, \texttt{\_mass} \\
    $S_{n},\ S_{t}$  & normal, tangential stiffness     & \unit{\newton\per\metre} & \texttt{s\_n}, \texttt{s\_t} \\
    $\boldsymbol{\xi}$ & tangential displacement        & \unit{\metre}        & \texttt{tangential\_spring} \\
    $\mathbf{F}_{n},\ \mathbf{F}_{t}$ & contact force components & \unit{\newton} & \texttt{f\_n}, \texttt{f\_t} \\
    $T_{i}$          & torque on particle $i$           & \unit{\newton\metre} & \texttt{torque[i]} \\
    $s_{p},\ v_{p}$  & platen translation, velocity     & \unit{\metre}, \unit{\metre\per\second} & \texttt{platen.offset}, \texttt{.velocity} \\
    $\mathbf{d}_{p}$ & platen direction of motion       & --                   & \texttt{platen.direction} \\
    $u_{p},\ F_{p}$  & platen displacement, reaction    & \unit{\metre}, \unit{\newton} & \texttt{platen.displacement}, \texttt{.reaction} \\
    $m_{p}$          & platen mass                      & \unit{\kilogram}     & \texttt{LoadedPlaten.mass} \\
    $F(t)$           & prescribed external load         & \unit{\newton}       & \texttt{LoadedPlaten.load} \\
    $\Delta t$       & time step                        & \unit{\second}       & \texttt{sim.dt} \\
    $T_{R}$          & Rayleigh critical time step      & \unit{\second}       & \texttt{rayleigh\_time\_step} \\
    $y_{w}$          & water table height               & \unit{\metre}        & \texttt{Water.level} \\
    $\rho_{w},\ \mu_{w}$ & water density and viscosity  & \unit{\kilogram\per\cubic\metre}, \unit{\pascal\second} & \texttt{density}, \texttt{viscosity} \\
    $U$              & superficial seepage velocity     & \unit{\metre\per\second} & \texttt{seepage\_velocity} \\
    $\varepsilon$    & local volume porosity            & --                   & \texttt{sim.porosity} \\
    $p$              & excess pore pressure             & \unit{\pascal}       & \texttt{excess\_pore\_pressure} \\
    $K_{w},\ k$      & water bulk modulus, permeability & \unit{\pascal}, \unit{\metre\squared} & \texttt{bulk\_modulus}, \texttt{permeability} \\
    \bottomrule
  \end{tabular}
\end{table}
